La réalisation de ce mémoire a bénéficié du soutien et de l'accompagnement de plusieurs personnes. Nous tenons à leur exprimer, par ces quelques lignes, notre profonde gratitude.

Nos remerciements vont particulièrement à l'endroit de :

\begin{itemize}[label=\ding{118}, leftmargin=1cm, itemsep=6pt]
    \item \textbf{Dr Belko Abdoul Aziz DIALLO}, notre encadreur, pour sa disponibilité, ses conseils et la rigueur avec laquelle il a suivi ce travail du début à la fin ;

    \item \textbf{Monsieur le Directeur de l'Institut Burkinabè des Arts et Métiers (IBAM)}, ainsi qu'à toute l'équipe pédagogique et administrative, pour la qualité de la formation reçue et leur accompagnement ;

    \item les chercheurs burkinabè dont les travaux sur les plantes médicinales constituent le corpus de cette étude, sans lesquels ce mémoire n'aurait pas pu voir le jour ;

    \item nos camarades de promotion et nos collègues, pour leur esprit d'entraide et les échanges qui ont enrichi ce travail ;

    \item notre famille, pour son soutien indéfectible depuis toujours, et particulièrement durant la préparation de ce mémoire ;

    \item tous nos amis et toutes les personnes qui, de près ou de loin, ont contribué à l'aboutissement de ce travail.
\end{itemize}
